\documentclass[a4paper]{article}

%% Language and font encodings
\usepackage[english]{babel}
\usepackage[utf8]{inputenc}
\usepackage[T1]{fontenc}

%% Page size and margins
\usepackage[a4paper,top=2cm,bottom=2cm,left=2.5cm,right=2.5cm,marginparwidth=1.75cm]{geometry}

%% Core packages
\usepackage{amsmath}
\usepackage{graphicx}
\usepackage{tabularx}
\usepackage{booktabs}
\usepackage{float}
\usepackage{subcaption}
\usepackage{microtype}
\usepackage{xcolor}
\definecolor{tcdBlue}{RGB}{5, 105, 185}

%% Hyperlinks
\usepackage[colorlinks=true,allcolors=.,urlcolor=blue]{hyperref}

%% Section heading style
\renewcommand{\familydefault}{\sfdefault}
\usepackage{titlesec}
\titleformat{\section}[hang]{\normalfont\Large\bfseries\color{tcdBlue}}{\thesection}{1em}{}{}

\hyphenation{op-tical net-works semi-conduc-tor}

\begin{document}

\title{\textbf{Representation Learning in Deep Reinforcement Learning:}\\
       \large A Comparative Study of DQN and Double DQN on Atari}
\author{Vinol Tauro\\
        \small Student Number: 25340566\\
        \small School of Computer Science and Statistics\\
        \small Trinity College Dublin}
\date{June 2026}
\maketitle

\begin{abstract}
This study investigates how structurally similar yet strategically distinct games shape the internal representations of deep reinforcement learning agents, how those representations respond to cross-game transfer, and what happens when an agent trained on one game is forced to learn another. Pong and Breakout share fundamental low-level visual structure but differ substantially in strategic demands. Using a 2$\times$2 factorial design, DQN and Double DQN (DDQN) agents were trained on both games and their 512-dimensional penultimate layer activations were analysed using t-SNE, Grad-CAM, dead neuron analysis, cosine similarity, and layer-wise Centred Kernel Alignment (CKA). The results reveal a representational hierarchy in which early convolutional layers remain highly similar across games (CKA $\geq$ 0.68) while the representation layer diverges almost completely (CKA $\leq$ 0.09), providing quantitative evidence that generalisation and specialisation are spatially separated within the network. Cross-game transfer experiments show that this geometric similarity does not translate to functional compatibility: zero-shot layer transplantation fails completely, and frozen conv backbones underperform scratch training. Sequential training experiments confirm that catastrophic forgetting is immediate and total, with Pong performance collapsing from $+9.1$ to $-21$ within 200,000 Breakout training steps. Crucially, the forgetting is localised to the representation layer rather than the convolutional layers: freezing the full backbone (conv and fc\_repr) prevents forgetting almost completely, while freezing convolutional layers alone is insufficient as the representation layer continues to adapt. These findings connect the representational hierarchy identified by CKA analysis to the mechanism of catastrophic forgetting, and establish layer-freezing depth as a direct lever for controlling the forgetting-plasticity tradeoff in continual learning.
\end{abstract}

\tableofcontents
\newpage

% ============================================================
\section{Introduction}
% ============================================================

Pong and Breakout look similar on the surface. Both feature a ball bouncing off a paddle, both share the same deflection physics, and both are rendered in the same pixelated Atari style. Yet anyone who has played both knows they demand completely different things from a player: Pong is about reading and reacting to an opponent, Breakout is about planning a path through a grid of bricks. The question this study asks is whether a deep reinforcement learning agent trained purely from pixels and reward signals picks up on that distinction internally, and if so, where the line falls between what is shared and what is game-specific in its learned representations.

This question has practical consequences beyond academic interest. As deep RL systems are deployed in more complex settings, understanding what agents learn internally, not just how well they score, becomes directly relevant to interpretability and transfer. Knowing which layers encode general visual features and which encode game-specific strategy has immediate implications for how knowledge can be reused or protected across tasks.

This study trains DQN \cite{mnih2015human} and Double DQN (DDQN) \cite{van2016deep} agents on both games and analyses the 512-dimensional activations at the penultimate fully connected layer. DQN and DDQN are used not because the algorithm comparison is the primary question, but because two independent algorithms provide two separate windows into the same representational phenomenon. If a finding holds across both, it is more likely to reflect something genuine about the games rather than an artefact of a particular training procedure.

The study proceeds in three phases. First, representation analysis establishes where cross-game similarity exists and at what layer it breaks down. Second, transfer experiments test whether that similarity translates to functional reusability. Third, sequential training experiments measure catastrophic forgetting and identify which layers are responsible. Each phase answers a question raised by the previous one, building toward a mechanistic account of how game-specific representations form and what happens when they are disrupted.

Four questions guide the work:

\begin{enumerate}
    \item Do agents trained on Pong and Breakout develop similar internal representations, and at which network depth does similarity break down?
    \item Does geometric similarity between layers translate to functional transferability across games?
    \item When an agent is trained sequentially on two games, how quickly and completely does it forget the first?
    \item At which layer does forgetting occur, and can it be prevented by freezing layers identified as general by CKA analysis?
\end{enumerate}

% ============================================================
\section{Background}
\label{sec:background}
% ============================================================

\subsection{Deep Q-Networks}

DQN \cite{mnih2015human} approximates the action-value function $Q(s, a)$ using a convolutional neural network. Two mechanisms stabilise training on raw pixel inputs: experience replay, which stores transitions in a buffer and samples them randomly to break temporal correlations; and a target network, a periodically frozen copy of the online network used to compute training targets. The DQN TD target is:

\begin{equation}
    y = r + \gamma \cdot \max_{a'} Q(s', a';\, \theta^{-})
\end{equation}

where $\theta^{-}$ are the target network parameters. Because the same network both selects and evaluates the best next action, taking the maximum over noisy Q-value estimates produces a systematic upward bias that compounds over training.

\subsection{Double DQN}

Van Hasselt et al.\ \cite{van2016deep} proposed using the online network to select the best next action and the target network to evaluate it:

\begin{equation}
    y = r + \gamma \cdot Q\!\left(s',\, \argmax_{a'} Q(s', a';\, \theta);\, \theta^{-}\right)
\end{equation}

Because the two networks have different error profiles, their errors are unlikely to compound in the same direction, substantially reducing overestimation bias. In this study's codebase, \texttt{DoubleDQNAgent} inherits everything from \texttt{DQNAgent} and overrides only the \texttt{learn()} method. Any representational difference between the two agents is attributable solely to this target formulation.

\subsection{Representation Learning and Capacity Loss}

The 512-dimensional activation vector at the penultimate fully connected layer serves as the agent's compressed encoding of the current game state. Zahavy et al.\ \cite{zahavy2016graying} showed that t-SNE projections of DQN representations reveal semantically meaningful clustering. Lyle et al.\ \cite{lyle2022understanding} demonstrated that deep RL agents progressively lose representational capacity as neurons enter the ReLU dead zone, a process accelerated by noisy gradient signals. This capacity loss is directly relevant to catastrophic forgetting: neurons that encoded features for one task and subsequently become dead under another task's gradient signal can never be reactivated.

\subsection{Catastrophic Forgetting and Continual Learning}

Catastrophic forgetting \cite{mccloskey1989catastrophic} refers to the abrupt loss of previously learned information when a neural network is trained on a new task. In deep RL, this manifests as performance collapse on the original game when training continues on a new one. The shared weights that encoded the first task's strategy are overwritten by the new task's gradient signal. Standard mitigations include Elastic Weight Consolidation \cite{kirkpatrick2017overcoming}, which penalises changes to weights important for the original task, and continual backpropagation \cite{elsayed2024continual}, which periodically reinitialises low-utility neurons to maintain plasticity. This study focuses on a simpler and more interpretable intervention: layer freezing, guided by the CKA analysis of which layers are general versus game-specific.

\subsection{Centred Kernel Alignment}

Centred Kernel Alignment (CKA) \cite{kornblith2019similarity} measures the similarity between two sets of representations, independent of linear transformations. Given activation matrices $X \in \mathbb{R}^{n \times p}$ and $Y \in \mathbb{R}^{n \times q}$ collected over the same $n$ inputs:

\begin{equation}
    \text{CKA}(X, Y) = \frac{\|Y^{\top}X\|_F^2}{\|X^{\top}X\|_F \cdot \|Y^{\top}Y\|_F}
\end{equation}

CKA returns values in $[0, 1]$, where 1 indicates identical representational geometry up to linear transformation. Both networks must be evaluated on the same inputs for CKA to be meaningful; using different game states per network would measure input dissimilarity rather than representational similarity.

\subsection{t-SNE and Grad-CAM}

t-SNE \cite{van2008visualizing} reduces high-dimensional data to two dimensions by minimising the KL divergence between Gaussian neighbourhood distributions in the original space and Student-t distributions in the 2D projection, faithfully preserving local cluster structure. Grad-CAM \cite{selvaraju2017grad} produces visual explanations by backpropagating gradients to the final convolutional layer and weighting feature maps by their importance to the output.

% ============================================================
\section{Experimental Design}
\label{sec:method}
% ============================================================

\subsection{Factorial Baseline}

The study uses a fully crossed 2$\times$2 factorial design, crossing algorithm (DQN, DDQN) with game (Pong, Breakout):

\begin{table}[H]
\centering
\caption{Factorial design. All four runs use identical hyperparameters and the same random seed.}
\label{tab:factorial}
\begin{tabular}{lcc}
\toprule
 & \textbf{Pong} & \textbf{Breakout} \\
\midrule
\textbf{DQN}  & Run 1 (2M steps) & Run 2 (5M steps) \\
\textbf{DDQN} & Run 3 (2M steps) & Run 4 (5M steps) \\
\bottomrule
\end{tabular}
\end{table}

\subsection{Preprocessing and Architecture}

Raw 210$\times$160 RGB frames are transformed through the standard DeepMind preprocessing stack \cite{mnih2015human}: no-op reset, frame skip (4), episodic life, grayscale and resize to 84$\times$84, reward clipping to $\{-1, 0, +1\}$, and frame stacking (4 frames).

Both agents share the CNN from Mnih et al.: three convolutional layers (32 filters, 8$\times$8, stride 4; 64 filters, 4$\times$4, stride 2; 64 filters, 3$\times$3, stride 1), followed by a 512-dimensional fully connected representation layer (\texttt{fc\_repr}) with ReLU activation, and a linear output head (\texttt{fc\_out}) producing Q-values for each action. A PyTorch forward hook on \texttt{fc\_repr} captures the 512-dim activation after every forward pass without modifying computation. Weights use orthogonal initialisation \cite{saxe2013exact}.

\begin{table}[H]
\centering
\caption{Hyperparameters used across all runs.}
\label{tab:hyperparams}
\begin{tabularx}{\linewidth}{lX}
\toprule
\textbf{Parameter} & \textbf{Value} \\
\midrule
Learning rate            & $1 \times 10^{-4}$ (Adam) \\
Replay buffer            & 100,000 transitions \\
Batch size               & 32 \\
Discount factor $\gamma$ & 0.99 \\
Target update frequency  & 1,000 steps (hard copy) \\
Epsilon schedule         & $1.0 \to 0.01$ over 100,000 steps \\
Gradient clipping        & $L_2$ norm $\leq 10.0$ \\
Checkpoint frequency     & Every 500,000 steps \\
Random seed              & 42 \\
\bottomrule
\end{tabularx}
\end{table}

\subsection{CKA Layer Analysis}

Forward hooks are registered at conv1, conv2, conv3, and fc\_repr. Activations are collected by running all four agents on 1,000 shared Pong frames obtained via random actions. Using random actions ensures a diverse, policy-neutral probe of the state space, and the shared frame set is essential for CKA validity: both networks must process the same inputs for the similarity measure to reflect representational geometry rather than input differences.

\subsection{Cross-Game Transfer Experiments}

\textbf{Zero-shot evaluation.} Chimera networks are constructed by transplanting convolutional layers from one trained agent into another game's upper layers, with no additional training. Ten conditions are evaluated across both games (50 episodes each), spanning native agents through cross-game and cross-algorithm chimeras to fully random networks.

\textbf{Backbone training.} DQN is initialised from the converged DQN/Pong checkpoint and trained on Breakout for 5M steps under three conditions: scratch (random initialisation baseline), frozen conv (Pong conv layers locked), and full fine-tune (all layers initialised from Pong, all trainable).

\subsection{Sequential Training and Forgetting Measurement}

Three conditions test catastrophic forgetting when a trained Pong agent continues training on Breakout for 2M steps:

\textbf{No freeze.} All layers adapt freely. Serves as the forgetting baseline.

\textbf{Freeze all.} Conv layers and fc\_repr are both loaded from Pong and frozen. Only a new Breakout output head (512$\times$4 = 2,052 parameters) trains. This tests whether freezing the full backbone prevents forgetting.

\textbf{Freeze conv.} Conv layers are frozen; fc\_repr and a new Breakout output head train. This tests whether freezing only the convolutional layers, as CKA suggests they are the most general, is sufficient to protect Pong performance.

Forgetting is measured every 200,000 Breakout training steps via three signals: (1) Pong reward, evaluated using a chimera network composed of the current backbone with the original Pong output head; (2) dead neuron fraction in fc\_repr; and (3) CKA drift, the similarity between current fc\_repr activations and the original Pong fc\_repr activations on a fixed set of Pong frames.

\subsection{Interleaved Training}

A single DQN agent with a shared backbone and two game-specific output heads alternates between Pong and Breakout episode by episode. The shared backbone (conv + fc\_repr) receives gradients from both games, allowing it to develop representations that serve both tasks simultaneously. Separate replay buffers are maintained per game. Training proceeds for 1M steps per game (2M total) from random initialisation, allowing a clean comparison with the sequential baseline.

% ============================================================
\section{Results}
\label{sec:results}
% ============================================================

\subsection{Training Performance}

\begin{table}[H]
\centering
\caption{Training results for all four baseline agents. Reward statistics over the final 50 episodes.}
\label{tab:results}
\begin{tabularx}{\linewidth}{lXXXX}
\toprule
\textbf{Agent} & \textbf{Episodes} & \textbf{Final Reward} & \textbf{Final Q} & \textbf{Peak Q} \\
\midrule
DQN / Pong      & 2,434   & $7.12 \pm 5.22$ & 2.64 & 2.95 \\
DDQN / Pong     & 2,457   & $7.50 \pm 5.03$ & 2.05 & 2.48 \\
DQN / Breakout  & 144,751 & $3.86 \pm 4.94$ & 4.06 & 5.39 \\
DDQN / Breakout & 144,948 & $4.00 \pm 5.04$ & 4.13 & 4.55 \\
\bottomrule
\end{tabularx}
\end{table}

Both algorithms reach comparable final reward on both games. On Pong, DQN achieves $7.12 \pm 5.22$ and DDQN $7.50 \pm 5.03$. On Breakout, DQN reaches $3.86 \pm 4.94$ and DDQN $4.00 \pm 5.04$. The near-identical task performance is important: any representational differences observed in subsequent analysis cannot be attributed to differences in what the agents learned to do, only to how they learned to represent it.

\begin{figure}[H]
    \centering
    \begin{subfigure}[b]{0.48\linewidth}
        \includegraphics[width=\linewidth]{figures/training_curves_pong.png}
        \caption{Pong.}
    \end{subfigure}
    \hfill
    \begin{subfigure}[b]{0.48\linewidth}
        \includegraphics[width=\linewidth]{figures/training_curves_breakout.png}
        \caption{Breakout.}
    \end{subfigure}
    \caption{Training curves (20-episode rolling average). Both algorithms reach comparable final reward.}
    \label{fig:training_curves}
\end{figure}

\subsection{Q-Value Overestimation}

DQN Q-values on Pong peak at 2.95 and settle at 2.64; DDQN peaks at 2.48 and settles at 2.05, approximately 29\% lower. On Breakout the gap is smaller, with DQN showing a higher peak (5.39 versus 4.55) consistent with overestimation being amplified during the most uncertain phases of early training. The cleaner gradient signal under DDQN provides the direct mechanism for the representational quality differences that follow.

\begin{figure}[H]
    \centering
    \includegraphics[width=0.7\linewidth]{figures/qvalue_overestimation.png}
    \caption{Mean maximum Q-value over training. DQN values drift upward; DDQN values remain lower and more stable.}
    \label{fig:qvalue}
\end{figure}

\subsection{t-SNE Representation Analysis}

The dominant structural feature of the t-SNE embeddings is separation by game. When all four agents are projected together, two primary clusters emerge corresponding to Pong and Breakout, with algorithm-level sub-structure visible within each (Figure~\ref{fig:all_agents}). The game played is the primary organising principle of the representation space; the algorithm used is secondary.

\begin{figure}[H]
    \centering
    \begin{subfigure}[b]{0.48\linewidth}
        \includegraphics[width=\linewidth]{figures/tsne_game_effect_dqn.png}
        \caption{DQN: Pong vs Breakout.}
    \end{subfigure}
    \hfill
    \begin{subfigure}[b]{0.48\linewidth}
        \includegraphics[width=\linewidth]{figures/tsne_game_effect_ddqn.png}
        \caption{DDQN: Pong vs Breakout.}
    \end{subfigure}
    \caption{Game effect: representations separate cleanly by game regardless of algorithm.}
    \label{fig:game_effect}
\end{figure}

Within each game, DDQN produces noticeably tighter cluster geometry with lower intra-cluster variance (Figure~\ref{fig:algo_effect}), consistent with a cleaner gradient signal producing more organised internal representations.

\begin{figure}[H]
    \centering
    \begin{subfigure}[b]{0.48\linewidth}
        \includegraphics[width=\linewidth]{figures/tsne_algo_effect_pong.png}
        \caption{Algorithm effect on Pong.}
    \end{subfigure}
    \hfill
    \begin{subfigure}[b]{0.48\linewidth}
        \includegraphics[width=\linewidth]{figures/tsne_algo_effect_breakout.png}
        \caption{Algorithm effect on Breakout.}
    \end{subfigure}
    \caption{Algorithm effect: DDQN produces tighter cluster geometry within each game.}
    \label{fig:algo_effect}
\end{figure}

\begin{figure}[H]
    \centering
    \includegraphics[width=0.6\linewidth]{figures/tsne_all_agents.png}
    \caption{All four agents. Game clusters dominate; algorithm sub-structure is visible within each.}
    \label{fig:all_agents}
\end{figure}

Representations coloured by episode reward show a partial spatial correspondence between t-SNE position and value, more consistently under DDQN (Figure~\ref{fig:reward}). Temporal evolution plots (Figure~\ref{fig:temporal}) show representations beginning as diffuse clouds and tightening progressively over training, with DDQN developing structure earlier on both games.

\begin{figure}[H]
    \centering
    \includegraphics[width=0.65\linewidth]{figures/tsne_by_reward.png}
    \caption{t-SNE coloured by cumulative reward. Spatial correspondence between value and position is more consistent under DDQN.}
    \label{fig:reward}
\end{figure}

\begin{figure}[H]
    \centering
    \begin{subfigure}[b]{0.48\linewidth}
        \includegraphics[width=\linewidth]{figures/tsne_temporal_dqn_pong.png}
        \caption{DQN / Pong.}
    \end{subfigure}
    \hfill
    \begin{subfigure}[b]{0.48\linewidth}
        \includegraphics[width=\linewidth]{figures/tsne_temporal_ddqn_pong.png}
        \caption{DDQN / Pong.}
    \end{subfigure}
    \vskip\baselineskip
    \begin{subfigure}[b]{0.48\linewidth}
        \includegraphics[width=\linewidth]{figures/tsne_temporal_dqn_breakout.png}
        \caption{DQN / Breakout.}
    \end{subfigure}
    \hfill
    \begin{subfigure}[b]{0.48\linewidth}
        \includegraphics[width=\linewidth]{figures/tsne_temporal_ddqn_breakout.png}
        \caption{DDQN / Breakout.}
    \end{subfigure}
    \caption{Temporal evolution across all four agents. Representations tighten over training; DDQN develops structure earlier.}
    \label{fig:temporal}
\end{figure}

\subsection{Saliency Maps}

Grad-CAM maps reveal consistent attention across all four agents: the ball attracts the highest saliency, with secondary attention on the agent's paddle. DDQN saliency is more spatially concentrated than DQN on both games, and shows greater forward-looking attention to bricks in the ball's trajectory on Breakout (Figures~\ref{fig:saliency_pong} and~\ref{fig:saliency_breakout}).

\begin{figure}[H]
    \centering
    \includegraphics[width=0.7\linewidth]{figures/saliency_pong.png}
    \caption{Grad-CAM saliency on Pong. Both agents attend to the ball; DDQN attention is more concentrated.}
    \label{fig:saliency_pong}
\end{figure}

\begin{figure}[H]
    \centering
    \includegraphics[width=0.7\linewidth]{figures/saliency_breakout.png}
    \caption{Grad-CAM saliency on Breakout. DDQN shows focused attention on the ball and bricks in its path.}
    \label{fig:saliency_breakout}
\end{figure}

\subsection{Dead Neurons and Cosine Similarity}

Dead neurons (those firing in fewer than 5\% of game states) are consistently higher under DQN than DDQN across both games (Figure~\ref{fig:dead_neurons}), consistent with the capacity loss findings of Lyle et al.\ \cite{lyle2022understanding}. Cross-game cosine similarity is non-trivial for both algorithms, with DDQN showing marginally higher alignment (Figure~\ref{fig:cosine}), suggesting its representations preserve more of the shared visual content.

\begin{figure}[H]
    \centering
    \includegraphics[width=0.7\linewidth]{figures/dead_neurons.png}
    \caption{Dead neuron fraction over training. DQN has consistently more inactive neurons than DDQN.}
    \label{fig:dead_neurons}
\end{figure}

\begin{figure}[H]
    \centering
    \includegraphics[width=0.7\linewidth]{figures/cosine_similarity.png}
    \caption{Cross-game cosine similarity. DDQN shows higher alignment, consistent with less gradient noise in its features.}
    \label{fig:cosine}
\end{figure}

\subsection{Layer-wise CKA Similarity}

CKA is computed between the Pong and Breakout agents at each of four layer depths using 1,000 shared Pong frames (Figure~\ref{fig:cka}). The shared frame set is essential: running each network on its own game's frames would measure input dissimilarity rather than representational similarity.

\begin{figure}[H]
    \centering
    \includegraphics[width=0.7\linewidth]{figures/layer_similarity_cka.png}
    \caption{Layer-wise CKA between Pong and Breakout agents. High similarity at early conv layers drops sharply at fc\_repr.}
    \label{fig:cka}
\end{figure}

Under DDQN, cross-game CKA is 0.92 at conv1, 0.85 at conv2, 0.41 at conv3, and just 0.09 at fc\_repr, representing a tenfold drop from early to late layers. The pattern holds under DQN (0.68, 0.72, 0.40, 0.05). Algorithm-effect comparisons within the same game show much higher similarity at all depths: on Pong, conv1 = 0.998, conv2 = 0.96, conv3 = 0.87, fc\_repr = 0.56, confirming that both algorithms develop functionally near-identical visual processing pipelines for the same game while differing in representational quality.

The sharpest drop in game-effect CKA occurs between conv3 and fc\_repr, falling from approximately 0.40 to 0.07. This marks the layer at which the network transitions from pixel-level feature extraction to state-value encoding, and where game-specific strategic demands take over from general visual processing.

\subsection{Cross-Game Transfer}

\subsubsection{Zero-Shot Evaluation}

\begin{figure}[H]
    \centering
    \begin{subfigure}[b]{0.48\linewidth}
        \includegraphics[width=\linewidth]{figures/mix_and_match_breakout.png}
        \caption{Evaluated on Breakout.}
    \end{subfigure}
    \hfill
    \begin{subfigure}[b]{0.48\linewidth}
        \includegraphics[width=\linewidth]{figures/mix_and_match_pong.png}
        \caption{Evaluated on Pong.}
    \end{subfigure}
    \caption{Zero-shot mix-and-match evaluation. Transplanting conv layers produces performance equivalent to random initialisation.}
    \label{fig:mixmatch}
\end{figure}

On Breakout, DQN/Pong conv transplanted into a DQN/Breakout agent achieves $0.44 \pm 0.78$, compared to $0.40 \pm 1.00$ for random conv and $0.42 \pm 0.80$ for a fully random network. The native agent scores $3.40 \pm 3.82$. On Pong, any foreign conv layer, whether from Breakout, a different algorithm, or random initialisation, produces a score of $-21$, the worst possible. High CKA at conv layers does not translate to functional compatibility: the fc\_repr layer was trained to interpret its own conv stack's specific activation patterns, and foreign conv activations fall outside that trained distribution.

\subsubsection{Backbone Training}

\begin{figure}[H]
    \centering
    \includegraphics[width=0.85\linewidth]{figures/backbone_comparison_reward.png}
    \caption{Breakout learning curves: scratch DQN vs frozen Pong conv vs full fine-tune from Pong. 500-episode rolling average.}
    \label{fig:backbone}
\end{figure}

\begin{table}[H]
\centering
\caption{Backbone training results on Breakout after 5M steps.}
\label{tab:backbone}
\begin{tabular}{lcc}
\toprule
\textbf{Condition} & \textbf{Final Reward} & \textbf{Best Smoothed} \\
\midrule
Scratch                   & 3.86 & 4.07 \\
Frozen conv (Pong)        & 2.28 & 3.03 \\
Full fine-tune from Pong  & 3.12 & 4.47 \\
\bottomrule
\end{tabular}
\end{table}

Frozen Pong conv underperforms scratch throughout training, despite CKA showing those layers are geometrically similar to what Breakout would develop independently. Pong's conv layers are tuned for features reward-predictive in Pong's visual context, such as opponent paddle trajectories and ball approach angles, which differ subtly from the features most useful for Breakout. Freezing them constrains rather than accelerates learning. Full fine-tune from Pong initialisation achieves the highest peak (4.47 smoothed), suggesting Pong weights provide a useful prior that accelerates early learning when all layers can subsequently adapt.

\subsection{Sequential Training: Catastrophic Forgetting}

Figure~\ref{fig:forgetting} shows the three forgetting metrics over 2M Breakout training steps for all three sequential conditions.

\begin{figure}[H]
    \centering
    \includegraphics[width=\linewidth]{figures/forgetting_curves.png}
    \caption{Forgetting curves across all three sequential conditions. Left: Pong reward. Centre: dead neuron fraction. Right: CKA drift from original Pong representations. The freeze-all condition maintains Pong performance throughout; freeze conv and no-freeze both show progressive degradation.}
    \label{fig:forgetting}
\end{figure}

\begin{table}[H]
\centering
\caption{Sequential training results. Pong reward, dead neurons, and CKA drift at start (step 0) and end (step 2M) of Breakout training.}
\label{tab:sequential}
\begin{tabularx}{\linewidth}{lXXXXXX}
\toprule
 & \multicolumn{2}{c}{\textbf{Pong Reward}} & \multicolumn{2}{c}{\textbf{Dead Neurons}} & \multicolumn{2}{c}{\textbf{CKA Drift}} \\
\cmidrule(lr){2-3}\cmidrule(lr){4-5}\cmidrule(lr){6-7}
\textbf{Condition} & Step 0 & Step 2M & Step 0 & Step 2M & Step 0 & Step 2M \\
\midrule
No freeze     & 9.1 & $-$20.8 & 0.879 & 0.873 & 0.037 & 0.041 \\
Freeze conv   & 9.1 & $-$14.0 & 0.879 & 0.836 & 0.055 & 0.043 \\
Freeze all    & 9.1 & \ \ \ 9.4  & 0.877 & 0.881 & 0.035 & 0.065 \\
\bottomrule
\end{tabularx}
\end{table}

\textbf{No freeze.} Pong reward collapses from $+9.1$ to $-20.75$ within the first 200,000 Breakout training steps and never recovers, reaching $-20.8$ at 2M steps. Dead neurons increase steadily from 0.729 at 200k steps to 0.867 at 2M, confirming progressive capacity loss as the mechanism underlying forgetting.

\textbf{Freeze conv.} Pong reward degrades more gradually but still substantially, falling to $-3.75$ by 200k steps and continuing to $-14.0$ at 2M steps. Dead neurons follow a similar trajectory. Freezing the convolutional layers alone is insufficient to prevent forgetting, because the representation layer fc\_repr continues to adapt freely to Breakout's strategic demands, progressively destroying the Pong-specific value encoding it had developed.

\textbf{Freeze all.} Pong reward remains stable throughout training, fluctuating between $7.2$ and $10.65$ across all 10 evaluation points and ending at $9.4$ after 2M Breakout steps. Dead neurons remain effectively unchanged at 0.875 to 0.881. Freezing both the convolutional layers and fc\_repr preserves the full Pong pathway and prevents forgetting almost completely.

The connection to the CKA analysis is direct. CKA identified fc\_repr as the layer with the sharpest cross-game divergence (CKA = 0.09). The sequential experiments confirm that this is also the layer through which forgetting propagates: when fc\_repr is free to adapt (no freeze, freeze conv), Pong performance collapses; when fc\_repr is frozen (freeze all), Pong performance is maintained. The forgetting is localised to exactly the layer that CKA predicted would need to specialise to game-specific strategic demands.

\subsection{Interleaved Training}

Interleaved training, in which a single agent with a shared backbone alternates episodes between Pong and Breakout from random initialisation, is ongoing at the time of writing. Results will be incorporated in the final version of this report. The experiment directly tests whether joint training on two structurally similar games prevents the forgetting observed in sequential training, and whether the shared backbone develops representations that generalise across both games as measured by cross-game CKA over the course of training.

% ============================================================
\section{Discussion}
\label{sec:discussion}
% ============================================================

\subsection{A Representational Hierarchy Confirmed at Multiple Levels}

The CKA depth profile establishes that early convolutional layers learn general visual features shared across games (CKA $>$ 0.68 at conv1 and conv2), while the representation layer specialises almost completely to game-specific strategic demands (CKA $<$ 0.09 at fc\_repr). The sequential training experiments provide independent causal confirmation of this hierarchy: forgetting propagates through fc\_repr, not through the convolutional layers. When fc\_repr is frozen, forgetting does not occur; when it is free to adapt, forgetting is immediate and total. The representational hierarchy identified by CKA is therefore not merely descriptive. It predicts precisely where intervention is needed to prevent forgetting.

\subsection{Geometric Similarity Does Not Imply Functional Compatibility}

High CKA at convolutional layers does not mean those layers can be reused without retraining. Zero-shot transplantation of conv layers fails completely, producing random-level performance, because fc\_repr was trained to interpret a specific conv stack's activation patterns. Foreign conv activations, even if geometrically similar, fall outside the trained distribution of fc\_repr and produce meaningless Q-value estimates. Similarly, frozen conv backbone training underperforms scratch, because Pong's conv features are tuned for Pong-specific visual statistics and constrain rather than support learning on Breakout. The lesson is that co-adaptation between adjacent layers is as important as the representational similarity of individual layers.

\subsection{The Forgetting-Plasticity Tradeoff}

The three sequential conditions reveal a clean spectrum. Freezing nothing produces complete forgetting but maximum plasticity for Breakout. Freezing conv only provides partial protection (Pong drops to $-14$ rather than $-21$), with the residual forgetting occurring through fc\_repr adaptation. Freezing the full backbone prevents forgetting but restricts Breakout learning to a single output head of 2,052 parameters. This spectrum represents a direct empirical measurement of the forgetting-plasticity tradeoff in a controlled RL setting. A practical agent operating under continual learning conditions would need to choose a point on this spectrum based on whether preserving old knowledge or acquiring new knowledge is the higher priority. The CKA analysis provides a principled guide for making that choice: layers with high cross-task CKA (conv1/2) are candidates for freezing; layers with low cross-task CKA (fc\_repr) drive both learning and forgetting.

\subsection{Algorithm Quality and Representation Stability}

DDQN consistently produces higher-quality representations: tighter t-SNE clusters, fewer dead neurons, more focused saliency, and higher cross-game cosine similarity. The root cause is the cleaner gradient signal from DDQN's decoupled target computation. This has implications for continual learning beyond this study: noisier training signals accelerate dead neuron accumulation, which is the neuron-level mechanism of capacity loss. An agent that retains more active neurons is better positioned to adapt to new tasks without catastrophic interference, since it has more representational capacity available for new encodings.

\subsection{Limitations}

\textbf{Single random seed.} All experiments use seed 42. Running multiple seeds (typically three to five in the literature \cite{henderson2018deep}) would be needed to confirm the observed patterns are robust rather than initialisation-specific. This is the most significant methodological limitation.

\textbf{Reduced replay buffer.} The 100,000-transition buffer is one-tenth of the 1M used in the original DQN paper \cite{mnih2015human}. Both algorithms use the same buffer, so the comparison is not confounded, but some effects may be more pronounced under the original conditions.

\textbf{Two games.} Results may not generalise beyond the specific structural relationship between Pong and Breakout. Replication across a broader set of game pairs with varying degrees of structural similarity would strengthen confidence in the findings.

\textbf{Layer granularity.} Freezing in this study operates at the layer level (conv group or fc\_repr). More fine-grained interventions, such as freezing individual feature maps identified as cross-game general by CKA, may offer better forgetting-plasticity tradeoffs than the coarse conditions tested here.

% ============================================================
\section{Conclusion}
\label{sec:conclusion}
% ============================================================

This study has built a mechanistic account of how game-specific representations form in deep RL agents, why they resist zero-shot transfer, and how they participate in catastrophic forgetting when training switches games. The key findings are:

\begin{enumerate}
    \item \textbf{Game content shapes representations more than algorithm.} t-SNE projections separate cleanly by game across both algorithms. Strategic demands, not visual similarity, drive representational divergence.

    \item \textbf{A representational hierarchy exists and can be localised.} CKA analysis reveals cross-game similarity of 0.92 at conv1 falling to 0.09 at fc\_repr under DDQN. The boundary between general visual processing and game-specific strategic encoding is at the conv-to-fc transition.

    \item \textbf{Geometric similarity does not imply functional compatibility.} Zero-shot conv transplantation fails completely; frozen conv backbone training underperforms scratch. Co-adaptation between adjacent layers determines functional transferability.

    \item \textbf{Catastrophic forgetting is immediate, total, and localised to fc\_repr.} Pong performance collapses from $+9.1$ to $-21$ within 200,000 Breakout training steps when all layers are free. Freezing conv alone is insufficient; forgetting propagates through fc\_repr. Freezing both conv and fc\_repr prevents forgetting almost completely.

    \item \textbf{The CKA hierarchy predicts the forgetting mechanism.} The layer identified as most game-specific by CKA (fc\_repr, CKA = 0.09) is exactly the layer through which forgetting propagates. CKA is not merely descriptive but predictive of where continual learning interventions need to act.

    \item \textbf{DDQN produces more stable representations.} Fewer dead neurons and tighter cluster geometry under DDQN suggest that cleaner gradient signals better preserve representational capacity, which has implications for continual learning robustness beyond this study.

    \item \textbf{Performance and representation quality are decoupled.} Both algorithms achieve near-identical task scores while producing representations of measurably different quality, confirming that reward alone is insufficient as an indicator of internal model quality.
\end{enumerate}

The central contribution of this work is the causal connection between the representational hierarchy and catastrophic forgetting: layers that CKA identifies as game-specific are the same layers through which forgetting propagates, and freezing those layers prevents forgetting. This connection provides a principled, analysis-guided approach to layer freezing in continual learning settings, one that does not require prior knowledge of which layers are general but can be determined from the CKA profile of independently trained agents.

% ============================================================
\newpage
\begin{thebibliography}{99}
\addcontentsline{toc}{section}{References}

\bibitem{mnih2015human}
Mnih V, Kavukcuoglu K, Silver D, Rusu AA, Veness J, Bellemare MG, et al.
Human-level control through deep reinforcement learning.
\textit{Nature}. 2015;518(7540):529--533.

\bibitem{van2016deep}
Van Hasselt H, Guez A, Silver D.
Deep reinforcement learning with double Q-learning.
In: \textit{Proceedings of the AAAI Conference on Artificial Intelligence}; 2016. vol.~30.

\bibitem{zahavy2016graying}
Zahavy T, Ben-Zrihem N, Mannor S.
Graying the black box: Understanding DQNs.
In: \textit{Proceedings of the 33rd International Conference on Machine Learning}; 2016. p.~1899--1908.

\bibitem{lyle2022understanding}
Lyle C, Rowland M, Dabney W.
Understanding and preventing capacity loss in reinforcement learning.
In: \textit{International Conference on Learning Representations}; 2022.

\bibitem{mccloskey1989catastrophic}
McCloskey M, Cohen NJ.
Catastrophic interference in connectionist networks: The sequential learning problem.
In: \textit{Psychology of Learning and Motivation}. 1989;24:109--165.

\bibitem{kirkpatrick2017overcoming}
Kirkpatrick J, Pascanu R, Rabinowitz N, Veness J, Desjardins G, Rusu AA, et al.
Overcoming catastrophic forgetting in neural networks.
\textit{Proceedings of the National Academy of Sciences}. 2017;114(13):3521--3526.

\bibitem{elsayed2024continual}
Elsayed M, Mahmood AR.
Addressing loss of plasticity and catastrophic forgetting in continual learning.
In: \textit{International Conference on Learning Representations}; 2024.

\bibitem{kornblith2019similarity}
Kornblith S, Norouzi M, Lee H, Hinton G.
Similarity of neural network representations revisited.
In: \textit{Proceedings of the 36th International Conference on Machine Learning}; 2019. p.~3519--3529.

\bibitem{van2008visualizing}
Van der Maaten L, Hinton G.
Visualizing data using t-SNE.
\textit{Journal of Machine Learning Research}. 2008;9(86):2579--2605.

\bibitem{selvaraju2017grad}
Selvaraju RR, Cogswell M, Das A, Vedantam R, Parikh D, Batra D.
Grad-CAM: Visual explanations from deep networks via gradient-based localization.
In: \textit{Proceedings of the IEEE International Conference on Computer Vision}; 2017. p.~618--626.

\bibitem{bellemare2013arcade}
Bellemare MG, Naddaf Y, Veness J, Bowling M.
The Arcade Learning Environment: An evaluation platform for general agents.
\textit{Journal of Artificial Intelligence Research}. 2013;47:253--279.

\bibitem{sutton2018reinforcement}
Sutton RS, Barto AG.
\textit{Reinforcement Learning: An Introduction}. 2nd~ed.
Cambridge, MA: MIT Press; 2018.

\bibitem{henderson2018deep}
Henderson P, Islam R, Bachman P, Pineau J, Precup D, Meger D.
Deep reinforcement learning that matters.
In: \textit{Proceedings of the AAAI Conference on Artificial Intelligence}; 2018. vol.~32.

\bibitem{saxe2013exact}
Saxe AM, McClelland JL, Ganguli S.
Exact solutions to the nonlinear dynamics of learning in deep linear neural networks.
In: \textit{International Conference on Learning Representations}; 2014.

\end{thebibliography}

\end{document}
